\section{The Common-Model Comparison Contract}\label{sec:model}
\subsection{Execution, resources and information}
CM-0.1 models one abstract device. A nonempty batch lasts exactly
$\Delta=0.001$ simulated seconds and may execute at most 32,768 work
tokens. Each prefill token costs one work unit; each scheduled decoding
request costs one unit and emits one token. A request's final prefill
batch emits no output, so decoding starts no earlier than the next batch.
Requests arriving during a batch become eligible at its end. When no work
is available, the clock advances directly to the next arrival.

Every resident reserves a full 32,768-token context, regardless of its
unknown output length. With capacity $K$, simultaneous residency is
bounded by $\min(64,\lfloor K/32768\rfloor)$. The two capacities,
131,072 and 262,144 token slots, therefore permit \emph{four and eight}
residents, respectively. Used KV is tracked separately from executed
prompt and output prefixes, rounded up to 16-token blocks. All frozen
prompt-plus-output lengths fit the context bound (maximum 30,053).
Reservation avoids output prediction and admission deadlock, at the cost
of conservative concurrency. It is not an incremental native KV allocator.

A policy sees only arrived request IDs and stable ordinals, arrival times,
prompt lengths, acknowledged progress, resident status, the last-output
or prefill-completion time, and the current clock. Actual or remaining
output length, source identity, future arrivals and output-derived
deadlines are hidden. The engine receives explicit per-request work spans
and either executes that plan exactly or rejects it. No eviction,
preemption, migration, prefix sharing or recomputation is available.
Completion releases the reservation. The common administrative horizon
is the last arrival plus 50 seconds; unfinished requests would remain
explicitly unfinished, not become observed latency samples.

\begin{table}[t]
\centering\small
\caption{Adaptations under CM-0.1. Names below denote project-maintained
rules, not executions of the original systems. All rows share the contract
in Section~\ref{sec:model}.}\label{tab:adaptations}
\begin{tabularx}{\linewidth}{lX}
\toprule Adaptation / lineage & Executed scheduling rule \\\midrule
Whole-FCFS / vLLM-inspired & Admit FCFS whole prompts when a slot and budget
permit; defer resident decodes in such batches. Otherwise decode residents. \\
Chunked-FCFS / chunked vLLM-inspired & Resident work before waiting work,
stable arrival order; prefill cap 512 tokens per request per batch. \\
Decode-first / Sarathi-inspired & Resident decodes, resident prefills,
then waiting prefills, each class FCFS; the same 512-token prefill cap. \\
SLO-gated / SLAI-inspired & Critical resident decodes first, then resident
prefills, waiting prefills by prompt length and ordinal, then remaining
decodes. Critical means time since last output/prefill completion at least
0.495 seconds. The prefill cap is 512 tokens. \\\bottomrule
\end{tabularx}
\end{table}

\subsection{Why some mechanisms coincide}
The four rules in Table~\ref{tab:adaptations} have distinct identities, but
need not induce distinct executions. At most eight residents or newly
admitted requests can receive work in a batch. With a 512-token prefill cap,
the three chunked rules use at most $8\times512=4096$ work tokens, well
below the 32,768 budget. Reordering resident decode and prefill work
therefore does not exclude any of it, and the fixed batch cost is
independent of order. Chunked-FCFS and decode-first admit waiting requests
in the same order and advance the same prefixes. Their equality is a
structural consequence in this domain, verified against all final records.

The same budget slack lets SLO-gated execute all eligible residual decodes
in every batch. The time since an acknowledged output or completed
prefill cannot accumulate to its 0.495-second gate. The gate is inactive
throughout these executions; prompt-length ordering of new admissions is
the relevant distinction. A gain attributed to an active urgency gate,
or to full SLAI, would misidentify the executed mechanism. Whole-FCFS can
differ because whole-prefill batches explicitly defer decodes and its
larger prompt spans can constrain admission within the budget.

\subsection{Conformance evidence}
An independent event referee reconstructs request progress, intent versus
executed spans, eligibility, token work, reservations, block-rounded KV,
timestamps, completion and metric accounting without importing the engine
or policy implementation. The execution gate also tested deterministic
replay, validated restart, input identity and future/latent-output
noninterference. All 5,760 journals passed the exhaustive in-allocation
check before final analysis. The subsequent review independently checked
all result and input identities, reran the separate summary verifier and
replayed 24 structurally selected journals. Those fresh sampled checks
supplement, rather than replace or purport to repeat, the exhaustive pass.
This evidence establishes conformance to the contract, not physical timing
accuracy or security isolation of policy code.
